%! Simplifying Radicals; Adding \& Subtracting — Unit 6, Lesson 6.2 — Homework
\documentclass[10pt]{article}
\usepackage{algebra2-article}
\usepackage{algebra2-boxes}
\newcommand{\TallMath}[1]{$\displaystyle #1\rule[-1.4em]{0pt}{3.2em}$}

\begin{document}

\pageheader{Unit 6, Lesson 6.2}{Homework}
\namedateperiod

\vspace{0.10in}

\begin{notesbox}{Practice}
Assume all variables represent \textbf{positive} real numbers. No calculator except where a problem
asks for a decimal.

\begin{enumerate}[leftmargin=1.9em, itemsep=11pt]
  \item \textbf{Simplify each numeric radical.} Watch the index --- it tells you the group size.
    \begin{center}
    \renewcommand{\arraystretch}{1.9}
    \begin{tabularx}{\linewidth}{@{}X X X X@{}}
    (a) \; $\sqrt{28}=$ \blank{1.7cm} & (b) \; $\sqrt{96}=$ \blank{1.7cm} &
    (c) \; $\sqrt{147}=$ \blank{1.7cm} & (d) \; $\sqrt{242}=$ \blank{1.7cm} \\
    (e) \; $\sqrt[3]{40}=$ \blank{1.7cm} & (f) \; $\sqrt[3]{-135}=$ \blank{1.7cm} &
    (g) \; $\sqrt[4]{80}=$ \blank{1.7cm} & (h) \; $\sqrt[5]{64}=$ \blank{1.7cm} \\
    \end{tabularx}
    \end{center}

  \item \textbf{Simplify each algebraic radical.} Divide the exponent by the index: quotient out,
        remainder in.
    \begin{center}
    \renewcommand{\arraystretch}{2.0}
    \begin{tabularx}{\linewidth}{@{}X X X@{}}
    (a) \; $\sqrt{x^{11}}=$ \blank{2.2cm} & (b) \; $\sqrt{20a^{3}}=$ \blank{2.2cm} &
    (c) \; $\sqrt{72m^{4}n^{5}}=$ \blank{2.4cm} \\
    (d) \; $\sqrt[3]{24p^{7}}=$ \blank{2.2cm} & (e) \; $\sqrt[3]{-16c^{9}}=$ \blank{2.2cm} &
    (f) \; $\sqrt{\dfrac{18y^{6}}{25}}=$ \blank{2.4cm} \\
    \end{tabularx}
    \end{center}

  \item \textbf{Add or subtract.} Simplify every term first. If a sum cannot be combined, say so ---
        but only after you have simplified.
    \begin{center}
    \renewcommand{\arraystretch}{2.0}
    \begin{tabularx}{\linewidth}{@{}X X@{}}
    (a) \; $6\sqrt{11}-2\sqrt{11}=$ \blank{2.6cm} &
    (b) \; $\sqrt{18}+\sqrt{32}=$ \blank{2.6cm} \\
    (c) \; $\sqrt{50}-\sqrt{8}+\sqrt{2}=$ \blank{2.6cm} &
    (d) \; $2\sqrt{63}+4\sqrt{28}=$ \blank{2.6cm} \\
    (e) \; $\sqrt{20}+\sqrt{12}=$ \blank{2.6cm} &
    (f) \; $\sqrt[3]{54}-\sqrt[3]{2}=$ \blank{2.6cm} \\
    (g) \; $\sqrt{45x}-\sqrt{20x}=$ \blank{2.6cm} &
    (h) \; $\sqrt{8}+\sqrt[3]{8}=$ \blank{2.6cm} \\
    \end{tabularx}
    \end{center}

  \item \textbf{Use the quotient property.}
    \begin{center}
    \renewcommand{\arraystretch}{2.0}
    \begin{tabularx}{\linewidth}{@{}X X X@{}}
    (a) \; $\sqrt{\dfrac{121}{9}}=$ \blank{2.0cm} &
    (b) \; $\dfrac{\sqrt{54}}{\sqrt{6}}=$ \blank{2.0cm} &
    (c) \; $\sqrt{\dfrac{5}{9}}=$ \blank{2.0cm} \\
    \end{tabularx}
    \end{center}

  \item \textbf{Two questions about what ``simplest'' means.}
    \begin{enumerate}[label=(\alph*), leftmargin=1.8em, itemsep=6pt]
      \item A classmate claims $\sqrt{36+64}=\sqrt{36}+\sqrt{64}$. Evaluate \emph{both} sides to
            show they are not equal, then write the rule that \emph{does} work and check it on the
            same two numbers. \writelines{3}
      \item Explain why $\sqrt[3]{16}$ is \emph{not} in simplest form but $\sqrt[3]{9}$ \emph{is} ---
            even though $9$ is a perfect square. \writelines{2}
    \end{enumerate}

  \item \textbf{Two roads, one answer (connecting to Lesson 6.1).} Simplify $\sqrt{x^{9}}$ twice.
    \begin{enumerate}[label=(\alph*), leftmargin=1.8em, itemsep=5pt]
      \item By splitting the radicand: $\sqrt{x^{9}}=\sqrt{x^{\blank{0.8cm}}\cdot x}=\blank{2.0cm}$
      \item With a rational exponent: $\sqrt{x^{9}}=x^{\blank{0.9cm}}=x^{\blank{0.9cm}}\cdot
            x^{\blank{0.9cm}}=\blank{2.0cm}$
      \item Why do these two methods always agree? Use the words \emph{mixed number} in your
            answer. \writelines{2}
    \end{enumerate}
\end{enumerate}
\end{notesbox}

\begin{extensionbox}
\textbf{The pendulum: quadruple the length to double the swing.} A pendulum of length $L$ feet takes
$T$ seconds to complete one full swing, where
\[
  T(L) = 2\pi\sqrt{\frac{L}{32}}\ \text{ seconds.}
\]
Every length below was chosen so that the fraction under the radical reduces to a perfect square.
\begin{center}
\renewcommand{\arraystretch}{1.3}
\begin{tabular}{|c|c|c|c|c|}
\hline
$L$ (feet)               & $2$ & $8$ & $18$ & $32$ \\ \hline
$\sqrt{L/32}$ (reduced)  & $\frac14$ & & & \\ \hline
$T(L)$ (exact)           & $\frac{\pi}{2}$ & & & \\ \hline
\end{tabular}
\end{center}
\begin{enumerate}[label=(\alph*), leftmargin=1.8em, itemsep=5pt]
  \item Complete the table \emph{exactly}, leaving $\pi$ in your answers. (Reduce the fraction
        \emph{before} taking the root --- simplest-form condition 2 says no fractions under a
        radical.)
  \item Give each period as a decimal, rounded to the nearest hundredth of a second.
  \item From $L=2$ to $L=8$ the length is multiplied by $4$. By what factor is the period
        multiplied? Check the same thing from $L=8$ to $L=32$. What does the exponent $\frac12$
        hidden in the square root do to the growth?
  \item A clockmaker wants a pendulum that swings exactly \emph{twice} as slowly as her current
        one. She reasons, ``I will just double the length.'' Why is she wrong, and what should she
        do instead?
  \item Show that $\sqrt{\dfrac{L}{32}}$ can also be written $\dfrac{\sqrt{2L}}{8}$, so that
        $T(L)=\dfrac{\pi\sqrt{2L}}{4}$. \emph{Hint:} multiply the top and bottom \emph{inside} the
        radical by $2$ and look at the new denominator. Verify your formula at $L=18$.
\end{enumerate}
\writelines{6}
\end{extensionbox}

\begin{spiralbox}
\textbf{Coming next --- Lesson 6.3: Multiplying \& Dividing Radicals; Rationalizing.} Today you ran
the product property in one direction, pulling perfect powers \emph{out}. Next you run it the other
way --- multiplying radicals together and distributing across sums, including FOIL on things like
$\left(3+\sqrt5\right)\left(2-\sqrt5\right)$. Then you finally clear simplest-form condition 3, the
one we skipped: a radical stuck in a \emph{denominator}. The tool for a binomial denominator is the
\textbf{conjugate} --- the same trick that turned $a+bi$ into a real number back in Lesson 3.4.
\end{spiralbox}

\end{document}
